\scanpage{4}
\mspage{3}
\diagram{\scriptsize\begin{tikzcd}[ampersand replacement=\&,column sep=.55em,row sep=1.3em]
{[\pi_2(Z)} \arrow[r] \& {\pi_2(Y)]} \arrow[r] \arrow[d] \& \pi_1(F) \arrow[r] \arrow[d,"\text{épi}"] \& \underline{\pi_1(Z)} \arrow[r] \arrow[d] \& \pi_1(Y) \arrow[r] \arrow[d,"\unsure{\sim|2|\wr}"] \& \pi_0(F) \arrow[r] \arrow[d,"\sim"] \& \underline{\pi_0(Z)} \arrow[r] \arrow[d] \& 1\\
\& \deleted{\unsure{*}}\ 1 \arrow[r] \& \pi_1(\Phi) \arrow[r] \& \pi_1(\widetilde Z) \arrow[r] \& \pi_1(\widetilde Y) \arrow[r] \arrow[d,phantom,"\wr"] \& \pi_0(\Phi) \arrow[r] \arrow[d,phantom,"\wr"] \& \pi_0(\widetilde Z) \arrow[r] \& 1\\
\& \& \& \& H \& \Sigma/G
\end{tikzcd}}
\diagram{\begin{tikzcd}[ampersand replacement=\&,row sep=tiny,column sep=1.3em]
F \arrow[d,hook] \& \deleted{\unsure{\Phi}} \& \Phi \arrow[d] \&\\
X \arrow[d] \& \deleted{\pi_1(X)=G} \arrow[d] \& \mathrm{Cat}\,G \arrow[d] \arrow[r,phantom,"\simeq"] \& \widetilde X \arrow[d]\\
S \& \deleted{\pi_1(S)} \& \mathrm{Cat}\,\Sigma \arrow[r,phantom,"="] \& \widetilde S
\end{tikzcd}}
$\Phi={}$ Catégorie des \deleted{éléments} de $\Sigma$, avec $\mathrm{Hom}_{\Phi}(\sigma,\sigma')={}$ ens des $g\in G$ tels que $\sigma=\sigma'f_1(g)$ i.e. $g\in f_1^{-1}(\sigma'^{-1}\sigma)$

Son $\pi_0$ est $\Sigma/G$, ses $\pi_1$ tous isom. à $\mathrm{Ker}\,f_1$.

On a $\Pi_1(F)\longrightarrow\Phi$, induisant une bijection sur les $\pi_0$, et sur les $\pi_1$ un épim., dont le noyau est $\mathrm{Coker}\,(\pi_2(X)\to\pi_2(S))$
\diagram{\scriptsize\begin{tikzcd}[ampersand replacement=\&,row sep=1.1em,column sep=.7em]
\pi_2(X) \arrow[r] \& \pi_2(S) \arrow[r] \& \pi_1(F) \arrow[r] \& \pi_1(X) \arrow[r] \arrow[d,phantom,"\wr"] \& \pi_1(S) \arrow[r] \arrow[d,phantom,"\wr"] \& \pi_0(F) \arrow[r] \arrow[d,phantom,"\wr"] \& 1\,\deleted{\unsure{\longrightarrow\pi_0}}\\
\& 1 \arrow[r] \& \pi_1(\widetilde F) \arrow[r] \& \pi_1(\widetilde X) \arrow[r] \& \pi_1(\widetilde S) \arrow[r] \& \pi_0(\Phi) \arrow[r] \& 1
\end{tikzcd}}

\[
\sigma=1\qquad \pi_1(\widetilde Z)=\{(g,h)\mid f_1(g)=g_1(h)
\]
$\pi_1(\widetilde Z)$ est ext. de $g_1^{-1}(\mathrm{Im}\,G)$ par $\pi_1(\Phi)$

$\pi_1(Z)$ est ext. de idem par $\mathrm{Coker}\,(\pi_2(Y)\to\pi_1(F))$

donc est ext. de $\pi_1(\widetilde Z)$ par
\[
\mathrm{Coker}\,(\pi_2(X)\times\pi_2(Y)\to\pi_2(S))
\]

\uline{Donc on trouve}

$Z\longrightarrow\widetilde Z$ induit un iso. pour les $\pi_1$ sss $\pi_2(X)\times\pi_2(Y)\longrightarrow\pi_2(S)$ est épi.

De plus, $\pi_i(Z)=0$ pour $i\geq 2$ \deleted{\unsure{si|\Leftrightarrow}} dès que $\pi_i(Y)=\pi_i(F)=0$ pour $i\unsure{=|\geq}2$, or $\pi_{i+1}(S)\to\pi_i(F)\to\pi_i(X)$ donc il suffit $\pi_i(X)=\pi_i(Y)=\pi_i(S)=0$ pour $i\geq 2$, i.e. $X,Y,S$ sont des $1$-tronqués\,\ldots
% archive sheet number: 3
